\documentclass[11pt]{article}
\usepackage[a4paper,margin=0.9in]{geometry}
\usepackage[colorlinks=true,linkcolor=blue]{hyperref}
\usepackage{amsmath,amssymb}
\usepackage{booktabs}
\usepackage{caption}
\usepackage{graphicx}
\usepackage{longtable}
\usepackage{titlesec}
\usepackage{xcolor}

% Lightweight callout boxes using only xcolor (no tcolorbox dependency,
% since this machine's BasicTeX install cannot be extended without sudo)
\newcommand{\intuitbox}[1]{\par\noindent\colorbox{blue!10}{\begin{minipage}{\dimexpr\textwidth-2\fboxsep\relax}\vspace{2pt}\textbf{Physical Intuition}\\[3pt]#1\vspace{2pt}\end{minipage}}\par\vspace{6pt}}
\newcommand{\codemapbox}[1]{\par\noindent\colorbox{green!10}{\begin{minipage}{\dimexpr\textwidth-2\fboxsep\relax}\vspace{2pt}\textbf{Where This Lives in the Code}\\[3pt]#1\vspace{2pt}\end{minipage}}\par\vspace{6pt}}

\title{The Array Geometry Handbook\\
\large How the Pinhole Array Design Code Works, and Why\\
\normalsize IYPT Pinhole Sunglasses Project --- A Student's Guide to the Physics and the Pipeline}
\author{}
\date{\today}

\begin{document}
\maketitle
\tableofcontents
\newpage

\section{Why This Folder Exists}

Imagine you have already solved the ``easy'' version of this problem: for a
single pinhole, what diameter gives the sharpest possible image for someone
with a given amount of short-sightedness? That question was answered earlier
in this project, and the answer changes with myopia power $M$ (in diopters)
--- a more short-sighted eye needs a smaller hole. That single-hole answer is
stored in \texttt{new\_sweeps/per\_M\_optimum\_map.csv} and is treated as a
\emph{fixed input} to everything in this handbook: we are not re-deriving hole
size here.

The harder question is: sunglasses need to block most of the light, and a
single tiny hole lets almost none through. The natural fix is to use
\emph{many} small holes instead of one. But the moment you have more than one
hole, a brand-new problem appears that a single hole never has: the eye can
form more than one image of the same scene at once, one through each hole.
If the eye is not perfectly focused (which, for someone wearing tinted
glasses rather than corrective lenses, is basically always true to some
degree), those images land in \emph{slightly different places} on the
retina. A small offset just looks like extra blur. A large offset looks like
literal double vision --- a real, disqualifying defect.

This handbook is about the code that answers one specific question:
\textbf{given a fixed hole size (already chosen to be optimal), how should
those many holes be arranged so that the double-vision problem never
happens, for every level of myopia and every direction the eye might be
looking?}

\subsection{The Two Physical Effects That Matter}

\intuitbox{Two physically independent things can go wrong with a multi-hole array, and
the code treats them as two separate checks, not one:
\begin{enumerate}
  \item \textbf{Ghosting.} If two holes are far enough apart, and the eye is
  out of focus, the two copies of the image they produce land in different
  places on the retina --- literally two overlapping pictures.
  \item \textbf{Vignetting.} As the eye rotates to look around a scene, the
  pupil --- a fixed-size circular opening --- physically slides across the
  pattern of holes. Holes that were far from the pupil edge at straight-ahead
  gaze can suddenly come into play once the eye turns, changing \emph{which}
  holes are simultaneously active and therefore changing the ghosting
  picture entirely.
\end{enumerate}
A pattern that is perfectly safe when the eye looks straight ahead can become
unsafe the moment the eye looks 15 degrees to one side. This is the single
most important finding in the whole study, and it is why the code always
checks both effects separately rather than assuming one implies the other.
}
\subsection{What the Code Does, in One Sentence}

For a chosen arrangement of holes (a \emph{geometry}), at a chosen hole
count and spacing, the code computes: (a) will any two holes ever produce a
visible double image, at any myopia level, and (b) will that stay true as
the eye looks around --- and if both answers are ``no'', the design is
\emph{structurally safe}. It then compares many such safe designs against
each other to find the one that packs the most holes (and therefore lets the
most light through) without crossing either safety line.

\newpage

\section{The Ghost-Separation Formula}

\subsection{Building the Intuition First}

Picture two pinholes, side by side, separated by a distance $s$ (their
center-to-center pitch). Light from a single point far away enters both
holes. If the eye's lens is perfectly focused for that light, both holes'
worth of light converges to \emph{exactly the same spot} on the retina ---
this is just what ``in focus'' means: parallel rays entering a lens at
different positions still meet at one point.

Now suppose the eye is short-sighted by $M$ diopters --- the lens converges
light slightly too much (or too little) for that point to land exactly on
the retina. The image plane is now displaced from the retina by a small
amount related to $M$. Each hole, because it sits at a different lateral
position, now projects its slightly-out-of-focus point to a \emph{slightly
different} spot on the retina. The two images are no longer perfectly
overlapping.

\intuitbox{The key physical fact that makes this predictable: the size of that
retinal offset grows in direct proportion to \emph{both} the hole
separation $s$ \emph{and} how far out of focus the eye is ($M$). Zero
separation, or zero defocus, means zero offset --- there is nothing subtle
here, it is basic similar-triangles geometry once you draw the ray paths
through the eye's lens.
}
\subsection{The Formula}

The validated closed-form result used throughout this project is:
\begin{equation}
\Delta r = \frac{s \cdot M}{P_0 - M}
\end{equation}
where $s$ is the hole-to-hole separation, $M$ is myopia in diopters, and
$P_0$ is the eye's total optical power (a fixed constant of the eye model,
$\approx 60$D). Notice the formula automatically gives $\Delta r = 0$ when
$M=0$ --- exactly the ``in focus, both holes agree'' case described above.
This is not a coincidence; it is the formula correctly reproducing a case we
already know the answer to by pure geometry, which is one of the strongest
kinds of sanity check available in physics.

An earlier, uncorrected version of this formula was used in this project
before it was properly checked. It did \emph{not} go to zero at $M=0$,
which on its own was enough to know it was wrong --- a formula that
disagrees with a case you can reason out by hand needs no further testing
to be rejected.

\subsection{How the Correct Formula Was Actually Confirmed}

Rather than trust either candidate formula on paper, this project built an
independent test: place exactly two holes in a real wave-optics simulation
(not a formula --- an actual propagation of light through the eye model),
measure where the two resulting bright spots land on the simulated retina,
and see which formula's prediction matches. Figure~\ref{fig:ghostvalid}
shows the result: the corrected formula fits the directly-measured data
with $R^2 = 0.995$ (essentially perfect agreement), while the old formula
scores $R^2 = -1701.7$ --- a negative $R^2$ means the old formula was
\emph{worse than simply guessing the average value every time}, which is
about as decisively wrong as a formula can be shown to be.

\begin{figure}[h]
\centering
\includegraphics[width=0.85\textwidth]{handbook_figs/ghost_formula_validation.png}
\caption{Directly-measured ghost separation (from real two-hole wave-optics
propagation) versus both candidate formulas. The corrected formula
($R^2{=}0.995$) is the one used throughout this project's array code; the
legacy formula ($R^2{=}-1701.7$) is not.}
\label{fig:ghostvalid}
\end{figure}

\codemapbox{The formula itself lives in \texttt{array\_geometry\_lab/degeneracy\_filter.py}
and \texttt{array\_optimization/stage0\_ghost\_validation/} (the original
validation, copied for speed into the geometry-comparison pipeline so it can
be evaluated thousands of times without re-running full wave optics each
time). The safety rule built on top of it is simple: for every pair of
holes that could be simultaneously lit by the pupil, compute $\Delta r$; if
the eye's own ability to resolve two separate points is smaller than every
$\Delta r$ in the pattern, the design is ghost-safe. This is
\texttt{gate\_ghost\_safety()} in \texttt{comparison\_metric/gates.py}.
}
\newpage

\section{The Vignetting Model}

\subsection{Why Ghosting Alone Is Not Enough}

The ghost formula above assumes we already know \emph{which pairs of holes}
are simultaneously letting light through to the eye. When the eye looks
straight ahead, the pupil sits centered behind the array, and it is easy to
list which holes fall under it. But the pupil is a physical opening of fixed
size, and the eye rotates constantly during normal viewing. As it rotates,
the pupil's position relative to the fixed hole pattern shifts --- some
holes that were fully under the pupil become partially or fully blocked,
and others that were outside the pupil's reach can come into play.

\intuitbox{This means a design that looks perfectly safe --- even \emph{trivially}
safe, with no two holes ever simultaneously active --- when evaluated only
at straight-ahead gaze can become unsafe the instant the eye looks even a
little to one side. You cannot conclude a design is safe from the
straight-ahead case alone; you have to check across the whole range of
gaze angles a person actually uses.
}
\subsection{The Geometry}

The pupil and a given hole are both circles. As the eye rotates by an angle
$\varphi$, the pupil's center shifts laterally relative to the hole pattern
by approximately
\begin{equation}
\Delta(\varphi) \approx R_{\rm rot} \sin\varphi
\end{equation}
where $R_{\rm rot}$ is the eye's physiological rotation radius
($\approx$13--14mm). The overlapping area between the pupil (radius
$R_p$) and a hole (radius $R_h$), once their centers are separated by
$\Delta$, is the standard geometric formula for the overlap of two circles:
\begin{equation}
A_{\rm overlap}(\Delta) = R_p^2\cos^{-1}\!\left(\frac{\Delta^2+R_p^2-R_h^2}{2\Delta R_p}\right)
+ R_h^2\cos^{-1}\!\left(\frac{\Delta^2+R_h^2-R_p^2}{2\Delta R_h}\right)
- \tfrac12\sqrt{(-\Delta{+}R_p{+}R_h)(\Delta{+}R_p{-}R_h)(\Delta{+}R_h{-}R_p)(\Delta{+}R_p{+}R_h)}
\end{equation}
with two edge cases handled separately: if the circles no longer touch at
all ($\Delta \ge R_p+R_h$), the overlap is exactly zero; if one circle sits
entirely inside the other ($\Delta \le |R_p-R_h|$), the overlap is simply
the area of the smaller circle. A na\"ive implementation of the general
formula can return a tiny negative number right at this second boundary,
purely from floating-point rounding --- a real bug this project found and
fixed by explicitly clamping the result to zero at the boundary.

\subsection{How the Implementation Was Checked}

The closed-form formula above was tested against a second, completely
independent method: randomly scatter thousands of points inside a hole's
area and simply count what fraction land inside the shifted pupil circle
(a Monte Carlo estimate). If the two methods --- one algebraic, one
brute-force random sampling --- agree, that is strong evidence the formula
was implemented correctly, since the two approaches share no code and could
not fail in the same way by coincidence. Figure~\ref{fig:vignvalid} shows
this comparison, converging to under 0.2 percentage points of relative
error.

\begin{figure}[h]
\centering
\includegraphics[width=0.85\textwidth]{handbook_figs/vignetting_validation.png}
\caption{Closed-form overlap-area formula versus an independent Monte Carlo
estimate, across gaze angle $\varphi$. Agreement here is what justifies
trusting the fast closed-form version for large sweeps instead of running
the (much slower) random-sampling method every time.}
\label{fig:vignvalid}
\end{figure}

\codemapbox{This geometry lives in \texttt{array\_optimization/vignetting.py}
(the original, validated implementation) and is copied into
\texttt{array\_geometry\_lab/degeneracy\_filter.py}\footnote{Copied rather
than imported for the large sweep pipeline, for speed --- but re-checked
byte-for-byte against the original's own validation data before use, so the
copy could not silently drift from the original.} for fast repeated use
across the geometry-comparison sweep. The safety rule built on it,
\texttt{gate\_vignetting\_safety()} in \texttt{comparison\_metric/gates.py},
re-runs the ``which holes are active'' calculation at many gaze angles
(0\textdegree\ to $\pm$25\textdegree) and requires the design to stay
ghost-safe at \emph{every} one of them, not just at $\varphi=0$.
}
\newpage

\section{The Nine Geometry Families}

With the two physical checks defined, the natural next question is: does the
\emph{shape} of the hole pattern matter? A square grid, a honeycomb, rings,
a spiral, or even a scattered random pattern all place $N$ holes over the
same area --- but they place them in very different relative positions,
which changes both the density you can achieve and the worst-case
hole-to-hole distances that feed the ghost formula.

\subsection{Deterministic Families}

Five families produce the \emph{same} layout every time, given the same
inputs (hole count, spacing) --- there is no randomness involved:
\begin{itemize}
  \item \textbf{Square} and \textbf{Hexagonal} lattices --- the simplest,
  most conventional packings. Hex packs more circles into a given area than
  square (about 90.7\% of the area versus 78.5\% at maximum density), which
  is why it tends to win on packing density when it survives the safety
  gates.
  \item \textbf{Rectangular} --- a square lattice allowed to have different
  spacing in the x and y directions, motivated by the fact that natural eye
  movement is not symmetric (people saccade horizontally more than
  vertically).
  \item \textbf{Concentric rings} --- holes arranged in circles around the
  center, a natural match to the pupil's own circular shape.
  \item \textbf{Fibonacci spiral / golden-angle phyllotaxis} --- the pattern
  seen in sunflower seed heads, where each new point is placed at a fixed
  angular step (the golden angle, $\approx$137.5\textdegree) from the last.
  These two names describe the mathematically identical pattern; the code
  confirms this explicitly rather than accidentally shipping the same
  generator twice under different names.
\end{itemize}

\subsection{Stochastic (Randomized) Families}

Two families place holes at \emph{different} positions every time they are
generated, even with identical inputs, because they use randomness:
\begin{itemize}
  \item \textbf{Poisson-disk (``blue-noise'')} --- random placement with a
  hard minimum-distance rule enforced between every pair of holes.
  \item \textbf{Jittered lattice} --- a hexagonal lattice with each hole's
  position nudged by a small random offset.
\end{itemize}
The reason to test these at all: every deterministic family above is
periodic (the pattern repeats at a fixed spacing), and periodic apertures
are known in optics to produce secondary bright spots at predictable angles
(diffraction grating orders) on top of whatever the ghost formula already
predicts. Breaking the periodicity is one way to avoid that side effect
entirely --- but a random pattern's safety now has to be checked
\emph{across many random realizations (seeds)}, not just once, since the
same requested parameters can produce a lucky safe layout on one seed and
an unsafe one on the next.

\begin{figure}[h]
\centering
\includegraphics[width=0.95\textwidth]{handbook_figs/section1_deterministic_catalog.png}
\caption{Representative generated layouts for the five deterministic
families, drawn to scale (hole radius shown as filled circles, not just
points) so packing density is visually honest.}
\end{figure}

\begin{figure}[h]
\centering
\includegraphics[width=0.95\textwidth]{handbook_figs/section1_stochastic_catalog.png}
\caption{Three different random seeds for the two stochastic families,
shown side by side. The visible differences between seeds are exactly why
these families need seed-averaged safety checking rather than a single
evaluation.}
\end{figure}

\codemapbox{Every geometry family is a small, self-contained function in
\texttt{array\_geometry\_lab/generators/}, one file per family, all sharing
one interface: given a requested hole count and spacing, return a list of
hole-center coordinates. Every generator also rejects (or removes) any hole
placement that would physically overlap its neighbors --- a fabrication
constraint separate from and prior to the optical safety checks. This
overlap-rejection can be aggressive: one documented test case showed a
requested 50 holes collapsing to only 9 actual holes once overlap was
enforced, at a spacing too tight for that hole size.
}
\newpage

\section{The Central Finding: Ghost-Safe Does Not Mean Vignetting-Safe}

This is the single most important experimental result to come out of the
whole geometry comparison, and it is worth stating plainly before looking
at any numbers: \textbf{a design can have an infinite ghost-safety margin at
straight-ahead gaze (meaning literally no two holes are ever simultaneously
active under the pupil) and still fail the vignetting check once the eye
looks around.}

\intuitbox{Why this makes physical sense: ``no two holes active at once'' at
$\varphi=0$ only describes one single eye position. As the pupil shifts
during rotation, holes that were previously outside its reach can come
into play together, creating an active pair that did not exist before.
Zero risk at one gaze angle says nothing about zero risk at every other
gaze angle --- these genuinely are two separate conditions to check.
}
Figure~\ref{fig:divergence} shows how often this actually happens, measured
directly from the data (not assumed): for every family, the fraction of
straight-ahead-safe designs that turn out to \emph{fail} the vignetting
check ranges from 37.8\% (concentric rings, the most robust family) to
88.9\% (Fibonacci spiral / golden angle, the least robust). This is not a
small-print caveat --- for the spiral families, roughly nine out of ten
designs that look perfectly safe under the simplest possible check turn out
to be unsafe once you account for how eyes actually move.

\begin{figure}[h]
\centering
\includegraphics[width=0.85\textwidth]{handbook_figs/fig_divergence.png}
\caption{Fraction of straight-ahead ghost-safe designs, per geometry family,
that fail the vignetting check once gaze angle is properly integrated.}
\label{fig:divergence}
\end{figure}

\codemapbox{This is why \texttt{comparison\_metric/gates.py} applies both gates as
independent, separately-required conditions
(\texttt{apply\_gates()} with both \texttt{gate\_ghost\_safety} and
\texttt{gate\_vignetting\_safety}), rather than treating one as a
shortcut proxy for the other.
}
\section{A Real Bug Worth Understanding: When ``Density'' Stops Making Sense}

\subsection{The Problem}

Packing density --- how many holes fit per unit area --- sounds like a
simple thing to compute: count the holes, divide by the area they occupy.
But ``the area they occupy'' is a genuinely tricky question once a
requested pattern collapses down to very few actual holes (because the
overlap-rejection step, described above, removed most of them).

\intuitbox{Picture a formula that assumes a repeating grid pattern (a ``unit cell''
multiplied by how many times it repeats). If overlap-rejection has reduced
a requested 20-hole pattern down to just \emph{one} surviving hole, there is
no repeating grid left to describe --- the formula's own assumption has
silently stopped applying, but nothing in the code was told to notice. The
formula still returns a number; it is simply a meaningless one.
}
This defect was real and it mattered: it was caught only because a
resulting design's reported density implied that its holes covered
\emph{more area than the space they were packed into} --- a value greater
than 100\% open area, which is a plain physical impossibility and an
excellent kind of bug to have a numerical assertion check for automatically,
rather than relying on a human noticing an oddly large number.

\subsection{The Fix, and the Lesson}

The fix required two separate corrections, because two different families
of geometry broke in two different ways:
\begin{enumerate}
  \item For lattice families (square, hex, rectangular), the fix is a
  \emph{feasibility check}: before trusting a unit-cell-based density, first
  confirm the requested spacing is even geometrically large enough to
  contain one hole. This check depends only on the ratio of hole size to
  spacing --- it is completely independent of how many holes actually
  survived, which is exactly why simply requiring ``at least 4 holes''
  \emph{could not} have caught this case on its own.
  \item For the ring/spiral/random families, the area was originally
  computed from the hole \emph{centers} only, forgetting that each hole
  also has its own physical radius extending past its center --- an
  omission proven with a real 7-hole test case where every hole was
  correctly non-overlapping, yet the naive area still came out too small.
  The fix adds each hole's own radius as a buffer around the boundary
  before computing area (a Minkowski/offset-buffer correction).
\end{enumerate}

The broader lesson, useful well beyond this one bug: \textbf{a formula that
is correct in the typical case can silently become meaningless in an edge
case, and the only reliable way to catch that is to check the answer
against a fact you know must always be true} --- here, that open area can
never exceed 100\%. Guessing at ``reasonable-looking'' thresholds (like a
minimum hole count) is not a substitute for that kind of physical sanity
check.

\codemapbox{The final, corrected logic lives in a single shared file,
\texttt{array\_geometry\_lab/degeneracy\_filter.py}, imported by every
downstream script rather than re-implemented separately in each one --- so
a future new geometry family only has to get this right once.
}
\newpage

\section{Choosing a Winner: The Two Comparison Metrics}

Once every candidate design has been checked against both safety gates,
many designs pass. The question becomes: among the safe ones, which is
actually best? ``Best'' needs a definition, and this project deliberately
tried two different ones rather than assuming a single answer.

\subsection{Metric 1: Maximum Packing Density}

Simple: among every design that clears both safety gates, pick the one
with the most holes per unit area. More holes per area means more open
area overall (for a fixed hole size), which means more light gets through
--- directly useful once transmission/brightness is brought back into the
picture later.

\intuitbox{Why not just rank by \emph{safety margin} instead? Because a design that is
safe with enormous margin (holes spaced very far apart) is not necessarily
useful --- it might only fit a handful of holes. Ranking by margin alone
would reward uselessly sparse designs. Density is the metric that actually
reflects the practical goal.
}
\subsection{Metric 2: The Pareto Front}

Metric 1 collapses everything to a single winning point, which hides any
trade-off that might exist between density and safety margin. Metric 2
instead reports the full \emph{Pareto front} --- every design for which no
other design is simultaneously denser \emph{and} safer. This shows the
whole trade-off curve rather than one point on it.

In the current dataset, this front collapses to a single point at every
myopia level --- meaning, at the tested spacing resolution, safety in this
project turned out to be close to binary (trivially safe versus clearly
unsafe) rather than a graded trade-off. This was checked to be a genuine
property of the data, not a bug in the Pareto logic itself, using a
synthetic test dataset built specifically to confirm the dominance-checking
code \emph{can} produce a multi-point front when one actually exists.

\codemapbox{Both metrics live in \texttt{comparison\_metric/metrics.py}, operating on
whatever survives \texttt{gates.py}'s filter --- deliberately kept as two
separate, swappable functions so a future third metric does not require
rewriting the comparison logic.
}
\section{Closing the Loop: Confirming the Theory Against Real Physics}

Everything above --- the ghost formula, the vignetting geometry, both
metrics --- uses fast, \emph{closed-form} formulas. They were validated once
against real wave-optics simulation for a simple two-hole case
(Section~2), but never checked against the specific, sometimes irregular,
multi-hole patterns that actually won the comparison. A formula validated
for a simple case is not automatically guaranteed to hold for a more
complex one built from the same physics.

\subsection{The Confirmation Step}

The winning designs from Metric~1 were each propagated through the
project's full multi-hole wave-optics engine (the same validated
simulation machinery used for the original single-aperture optimization,
not a shortcut) --- and the resulting simulated image was inspected
directly, with actual peak-detection on the simulated retina, rather than
trusting the formula's prediction.

\intuitbox{This distinction matters: the closed-form formula \emph{predicts} where
ghost images should land. The full simulation \emph{measures} where they
actually land, by literally propagating light through the eye model and
looking at the resulting image. Confirmation means these two independent
methods agree; it is not simply re-running the same formula a second time.
}
The result: \textbf{9 out of 9 finalist designs were confirmed} --- the
full simulation reproduced each finalist's predicted ghost-safety and
vignetting-safety exactly, at every tested gaze angle. One genuinely useful
side-finding also came out of this step: three of the nine finalists
revealed that an older reference file (\texttt{blur\_sweep.csv}) had been
computed at a since-corrected, lower simulation resolution --- caught
because three finalists with \emph{different} hole sizes reported the
\emph{identical} stale blur value to many decimal places, which is not
something real, correctly-computed diffraction physics would ever do for
different apertures. That reference file has since been regenerated at
the corrected resolution.

\codemapbox{This confirmation pass lives in
\texttt{array\_geometry\_lab\_phase2\_confirmation/}, importing (not
copying) the original validated simulation engine directly --- correctness
matters more than speed here, since only a handful of finalist designs
need checking, unlike the thousand-plus-candidate sweep that justified
using fast closed-form copies everywhere else.
}
\section{Summary: What a Design Needs to Pass}

\begin{center}
\begin{tabular}{cl}
\toprule
\textbf{Step} & \textbf{Requirement} \\
\midrule
1 & Hole size fixed by the existing acuity-optimal chart, per myopia level \\
2 & Generate a candidate geometry (pattern, hole count, spacing) \\
3 & \textbf{Gate 1:} no ghost-image pair violates the resolvable-separation limit \\
4 & \textbf{Gate 2:} Gate 1 holds at every tested gaze angle, not just $\varphi=0$ \\
5 & Among survivors, rank by packing density (Metric 1) or inspect the trade-off front (Metric 2) \\
6 & Confirm the winner(s) against full wave-optics simulation, not just the formula \\
\bottomrule
\end{tabular}
\end{center}

A design that passes all six steps has been checked, independently, at two
different levels of physical detail (fast formula and full simulation), and
is the point at which this project's theoretical work hands off to physical
fabrication and measurement.

\end{document}

